\scanpage{72}
On trouve que $J=\{\unsure{\delta^B_{z,i}|J^B_{z,i}|\xi^B_{z,i}}\}_{i\in I-\{z\}}$ -- en \unsure{termes|types} de base $(\delta_i)_{1\leq i\unsure{<|\leq}6}$, $\unsure{\dot z|z}=(\delta_6,\delta'_6)$, $J=\{\delta_{i6}\}_{1\leq i\leq5}$. La cat. des \deleted{$\frac12$ étoiles} couples $(C,J)$, $J$ une $\frac12$ étoile, est donc équivalente à celle des couples $(\tau,J)$, $\tau$ ens de card $2$, $J$ ens de card. $5$.
\margin{$\unsure{\widehat{\mathrm{Ens}}_2|\mathrm{Ens}_2}\times\widehat{\mathrm{Ens}}_5$}

\uline{Bases et bibases} (Il fallait commencer par \unsure{ça|là}!)
\margin{$\widehat{\mathrm{Ens}}_6$}
\margin{$\widehat{\mathrm{Ens}}_2\times\widehat{\mathrm{Ens}}_6$}
% slanted
$W=\mathrm{Aut}(C)$ est simplement transitif sur l'ens des bases épinglées ($=$ ordonnées). Il y a $72$ bases, le groupe de stabilité d'une base est $\simeq\unsure{\mathfrak S|\mathscr S}_6$. Il y a $36$ bibases, groupe de stabilité $\mathfrak S_2\times\mathfrak S_2$.
La cat. des $(C,b)$ est $\simeq$ celle des ens. de card $6$, celle des $(C,B=\{b,b'\})$ à celle \unsure{des|ses} couples $(\tau,I)$, $\tau$ ens de cardinal $2$, $I$ ens de cardinal $6$.
% archive sheet number: 71
